\chapter{Convolutions For Sequences}
\label{ch:convolutions-for-sequences}

\section{Why Start With Convolutions?}

Modern state space models are often introduced as alternatives to RNNs and transformers,
but the cleanest entry point is convolution. S4, S5, and related models can all be read as
ways to construct long, structured sequence filters. Before worrying about HiPPO matrices,
diagonal state transitions, or selective scans, we should be comfortable with this simpler
idea:
\[
  \text{a sequence model can be a learned filter applied across time.}
\]

You already know the transformer view: each token representation is updated by mixing
information from other positions. You also know the RNN view: a hidden state carries a
summary forward one step at a time. The convolutional view is different. It says that the
output at time $t$ can be computed from a fixed local pattern around $t$, using the same
weights at every time. This weight sharing is what makes convolution useful, and it is
also what gives us the first bridge to state space models.

The CNN vocabulary in this chapter is standard; see \citet{goodfellow2016deep} and the
CS231n notes \citep{cs231n} for the broader image-modeling treatment. Here we immediately
specialize it to one-dimensional temporal sequences, because that is the setting needed
for deep SSMs and for neural time series.

\section{Sequences, Channels, And Shapes}

Let an input sequence be
\[
  \vu_1,\vu_2,\ldots,\vu_L,
  \qquad
  \vu_t \in \R^{C_{\mathrm{in}}},
\]
where $L$ is the sequence length and $C_{\mathrm{in}}$ is the number of input channels.
For neural data, a channel might be a spike-count feature, a threshold-crossing feature,
or a spike-band-power feature. For language, it might be one coordinate of an embedding.
For audio, it might be one feature in a spectrogram frame.

A one-dimensional convolutional layer maps this sequence to
\[
  \vy_1,\vy_2,\ldots,\vy_{L'},
  \qquad
  \vy_t \in \R^{C_{\mathrm{out}}},
\]
where $C_{\mathrm{out}}$ is the number of output channels and $L'$ depends on padding,
stride, and kernel size. In deep-learning libraries, a batch of temporal data is often
stored as one of the following:
\[
  (B, L, C)
  \quad \text{or} \quad
  (B, C, L),
\]
where $B$ is batch size. The math is the same either way, but being explicit about the
time axis prevents many implementation mistakes.

\section{Convolution Versus Cross-Correlation}

Textbooks define discrete convolution between scalar sequences $u$ and $k$ as
\[
  (k \conv u)_t = \sum_i k_i u_{t-i}.
\]
The index $t-i$ means the kernel is flipped relative to the input. Many deep-learning
libraries actually implement cross-correlation:
\[
  (k \star u)_t = \sum_i k_i u_{t+i}.
\]
The distinction matters in signal processing, but it is usually not important for learned
CNNs because the kernel values are trainable. If the model can learn $k_0,k_1,k_2$, it can
also learn the reversed ordering. In these notes, we will use the causal sequence-modeling
convention
\begin{equation}
  \vy_t = \sum_{i=0}^{K-1} \mK_i \vu_{t-i}.
  \label{eq:causal-conv-basic}
\end{equation}
This is a true past-looking filter: $\vy_t$ depends on $\vu_t,\vu_{t-1},\ldots,\vu_{t-K+1}$
and not on future inputs.

The kernel slice $\mK_i$ is a matrix
\[
  \mK_i \in \R^{C_{\mathrm{out}} \times C_{\mathrm{in}}},
\]
so \eqref{eq:causal-conv-basic} means
\[
  \vy_t[c_{\mathrm{out}}]
  =
  \sum_{i=0}^{K-1}
  \sum_{c_{\mathrm{in}}=1}^{C_{\mathrm{in}}}
  \mK_i[c_{\mathrm{out}},c_{\mathrm{in}}]\,
  \vu_{t-i}[c_{\mathrm{in}}].
\]
This is the multi-channel version of a finite impulse-response filter. The layer uses
the same kernel matrices at every time $t$, which is the temporal weight-sharing
assumption.

\section{Causality}

For offline image classification, a convolution can use pixels on both sides of a location.
For a causal decoder, that would be future leakage. A causal model must produce the output
at time $t$ using only information available at or before $t$:
\[
  \vy_t = f(\vu_1,\ldots,\vu_t).
\]
Equation \eqref{eq:causal-conv-basic} satisfies this by construction. A centered
convolution such as
\[
  \vy_t = \mK_{-1}\vu_{t-1} + \mK_0\vu_t + \mK_{1}\vu_{t+1}
\]
does not. It may be perfectly appropriate for offline smoothing or bidirectional
representation learning, but it is not causal.

This distinction will stay important throughout the notes. S5 and Mamba can both be run
causally. Some training objectives or architectural wrappers can still accidentally leak
future information if masking, patching, padding, or normalization is handled carelessly.

\section{Padding And Boundary Effects}

The formula
\[
  \vy_t = \sum_{i=0}^{K-1} \mK_i \vu_{t-i}
\]
refers to inputs such as $\vu_0,\vu_{-1}$ near the beginning of the sequence. A padding
rule defines what happens there.

\paragraph{Valid convolution.}
Only compute outputs when the full window exists. For a causal kernel of length $K$, this
starts at $t=K$. There is no artificial boundary value, but the output sequence is shorter.

\paragraph{Zero padding.}
Set out-of-range inputs to zero. This preserves length, but early outputs see many zeros.
That is often fine, but the model can learn boundary artifacts if the sequence start has
special meaning.

\paragraph{Replication or reflection padding.}
Use repeated or reflected boundary values. These are common in image processing, less
central for causal sequence modeling, and often inappropriate for online decoding because
reflection can use future values.

For causal sequence models, the usual length-preserving choice is left zero padding:
\[
  \vu_t = 0 \quad \text{for } t \leq 0.
\]
Then $\vy_1$ is computed from mostly padding and $\vu_1$, while later positions see a full
history window. In long sequences the boundary region is usually negligible, but in short
segments or patch-based training it can be a real source of distribution shift.

\section{Stride And Downsampling}

Stride means the layer does not produce an output at every input time. With stride $s$,
one common causal convention is
\[
  \vy_j = \sum_{i=0}^{K-1} \mK_i \vu_{js-i}.
\]
If $s=2$, the layer emits one output for every two input steps. Stride is therefore both
a computational choice and a modeling choice: it reduces temporal resolution.

For image CNNs, stride is often a spatial downsampling tool. For temporal neural data, it
also changes the time scale of the representation. If bins are $20$ ms wide, stride $5$
means each output step advances by $100$ ms. That may be useful for some high-level
features, but it can be too coarse for precise speech timing.

Stride should not be confused with kernel size. A kernel of size $K=5$ with stride $1$
looks at five time bins and advances by one bin. A kernel of size $K=5$ with stride $5$
looks at five time bins and advances by five bins. Both have the same local window, but
the second throws away intermediate output positions.

\section{Receptive Field}

The receptive field of an output position is the set of input positions that can affect
it. For a single causal convolution with kernel size $K$ and stride $1$, the receptive
field of $\vy_t$ is
\[
  \{\vu_t,\vu_{t-1},\ldots,\vu_{t-K+1}\},
\]
so its size is $K$.

Stacking convolutional layers increases the receptive field. Suppose we stack $N$ causal
convolutions, each with kernel size $K$ and stride $1$. Then the final output at time $t$
depends on
\[
  1 + N(K-1)
\]
input positions. For example, five layers of kernel size $3$ see
\[
  1 + 5(3-1) = 11
\]
positions. This linear growth is a limitation if we need long context.

This is one reason long-sequence models become interesting. A transformer can access all
previous tokens in one layer under causal attention, but with quadratic attention cost in
sequence length. A plain CNN is cheap and parallel, but a shallow one sees only local
context. Deep SSMs try to get long effective memory with linear or near-linear cost.

\section{Dilated Convolution}

Dilated convolution expands receptive field without increasing kernel size. A causal
dilated convolution with dilation $d$ is
\[
  \vy_t = \sum_{i=0}^{K-1} \mK_i \vu_{t-di}.
\]
If $d=1$, this is the ordinary causal convolution. If $d=4$ and $K=3$, the output uses
$\vu_t$, $\vu_{t-4}$, and $\vu_{t-8}$.

Stacking layers with increasing dilation can grow the receptive field quickly. A common
pattern is
\[
  d = 1,2,4,8,\ldots.
\]
This gives a large context using few layers, and is one reason dilated temporal CNNs have
been successful in audio and sequence modeling.

The tradeoff is that dilation samples the past on a sparse grid inside each layer. Later
layers can combine information across these grids, but the architecture still imposes a
particular multiscale pattern. State space models will take a different route: they use a
recurrence to generate a long dense filter, often with smoothly decaying and oscillatory
modes.

\section{Channel Mixing}

The scalar convolution formula hides an important issue: how do features interact across
channels? In a standard multi-channel convolution, every output channel can use every
input channel at every kernel offset:
\[
  \mK_i \in \R^{C_{\mathrm{out}} \times C_{\mathrm{in}}}.
\]
The number of parameters is
\[
  K C_{\mathrm{in}} C_{\mathrm{out}}.
\]
This is expressive, but expensive when both channel counts are large.

\subsection{Pointwise Convolution}

A pointwise convolution is a kernel of size $1$:
\[
  \vy_t = \mW \vu_t,
  \qquad
  \mW \in \R^{C_{\mathrm{out}} \times C_{\mathrm{in}}}.
\]
It mixes channels but does not mix time. In sequence models, pointwise layers are often
used before or after a temporal operation. For example, a block might project from model
width $H$ to an expanded hidden width, run a temporal operation, then project back.

\subsection{Depthwise Convolution}

A depthwise convolution applies a separate temporal filter to each channel:
\[
  y_t[c] = \sum_{i=0}^{K-1} k_i[c]\,u_{t-i}[c].
\]
There is temporal mixing within each channel, but no channel mixing. The parameter count
is only $K C_{\mathrm{in}}$ if $C_{\mathrm{out}}=C_{\mathrm{in}}$.

Depthwise convolution is often paired with pointwise convolution:
\[
  \text{channel mixing} \rightarrow \text{depthwise temporal mixing} \rightarrow
  \text{channel mixing}.
\]
This factorization is less expressive than a full convolution but much cheaper. Mamba
uses a small local convolution inside its block, and this kind of separation between
local temporal processing and channel projection is a recurring design pattern.

\subsection{Grouped Convolution}

Grouped convolution lies between full and depthwise convolution. Channels are partitioned
into groups, and convolution is full within each group but not across groups. This reduces
parameter count and computation while allowing more channel interaction than a purely
depthwise operation.

\section{Convolution As A Matrix}

For a finite sequence, convolution can be written as multiplication by a structured
matrix. In the scalar causal case with kernel $(k_0,k_1,k_2)$ and left zero padding,
\[
\begin{bmatrix}
y_1 \\
y_2 \\
y_3 \\
y_4
\end{bmatrix}
=
\begin{bmatrix}
k_0 & 0   & 0   & 0 \\
k_1 & k_0 & 0   & 0 \\
k_2 & k_1 & k_0 & 0 \\
0   & k_2 & k_1 & k_0
\end{bmatrix}
\begin{bmatrix}
u_1 \\
u_2 \\
u_3 \\
u_4
\end{bmatrix}.
\]
This is a lower-triangular Toeplitz matrix: each diagonal is constant. The lower-triangular
structure is causality. The constant diagonals are time invariance.

This matrix view is useful for three reasons.

\begin{enumerate}
  \item It makes convolution look like a linear layer with heavy weight sharing.
  \item It shows why convolution is parallel over time: all outputs are part of one
  structured matrix-vector product.
  \item It prepares us for SSMs, where the same kind of Toeplitz convolution matrix can
  be generated from recurrence parameters rather than learned directly.
\end{enumerate}

Of course, deep-learning libraries do not literally build this large matrix for ordinary
small convolutions. They use specialized kernels. But mathematically, the structured
matrix picture is accurate.

\section{Finite Impulse Response Filters}

A causal convolution with kernel length $K$ is a finite impulse-response filter, usually
abbreviated FIR. To see this, imagine a scalar input sequence that is zero everywhere
except at time $1$:
\[
  u_1 = 1,
  \qquad
  u_t = 0 \text{ for } t \neq 1.
\]
For a causal convolution,
\[
  y_1 = k_0,\quad
  y_2 = k_1,\quad
  \ldots,\quad
  y_K = k_{K-1},
\]
and then the response becomes zero. The kernel is exactly the impulse response.

This language matters because linear state space models produce impulse responses too.
The difference is that an SSM impulse response can be very long, even infinite in
principle, while being parameterized compactly by state matrices.

\section{Convolution And Parallelism}

A vanilla RNN computes
\[
  \vh_t = f(\vh_{t-1}, \vu_t).
\]
This creates a sequential dependency: to compute $\vh_t$, we first need $\vh_{t-1}$.
That does not mean RNNs cannot be optimized, but it does limit straightforward parallelism
over time.

A convolution computes
\[
  \vy_t = \sum_{i=0}^{K-1} \mK_i \vu_{t-i}.
\]
Once the input sequence and kernel are known, every $\vy_t$ can be computed independently
of every other $\vy_{t'}$. The outputs share weights, but they do not depend on one
another. This is why CNNs train efficiently on long sequences.

The catch is memory length. A small convolution has a short receptive field. A large
convolution can have a long receptive field, but a directly learned dense kernel of length
$L$ has many parameters and can be expensive. The S4/S5 idea is to get the benefits of
long convolution while using a structured parameterization inherited from dynamical
systems.

\section{FFT Convolution}

There is another reason convolution is computationally special: long convolutions can be
accelerated using the Fourier transform. The convolution theorem says, informally,
\[
  \mathcal{F}(k \conv u) = \mathcal{F}(k) \odot \mathcal{F}(u),
\]
where $\odot$ is elementwise multiplication. So a long convolution can be computed by:
\[
  \text{FFT} \rightarrow \text{elementwise multiply} \rightarrow \text{inverse FFT}.
\]
For very long kernels, this can be much cheaper than directly summing over every offset.

This is one reason the convolutional view of SSMs is valuable. If an SSM can produce a
long kernel $\mK_0,\ldots,\mK_{L-1}$, then training can sometimes use fast convolution
methods instead of stepping the recurrence one time step at a time. S4 famously leans on
this distinction: recurrent for online generation, convolutional for parallel training.

There are details we are postponing. FFT convolution requires care with padding because
the discrete Fourier transform naturally computes circular convolution. To get ordinary
linear convolution, one pads the sequences before transforming. The key idea for now is
only that convolution has global algorithms that exploit structure.

\section{Convolutional Sequence Blocks}

A single linear convolution is usually not enough. Neural sequence models use convolution
inside nonlinear blocks. A simple temporal CNN block might look like
\[
  \vu_{1:L}
  \rightarrow
  \text{LayerNorm}
  \rightarrow
  \text{causal convolution}
  \rightarrow
  \text{nonlinearity}
  \rightarrow
  \text{pointwise projection}
  \rightarrow
  \text{residual addition}.
\]
The exact normalization and activation choices vary. The important point is that the
temporal operation is embedded in a deep residual architecture, just as attention is
embedded in transformer blocks and SSM layers are embedded in S4/S5/Mamba blocks.

A residual block typically has the schematic form
\[
  \vz_{1:L}^{(n+1)}
  =
  \vz_{1:L}^{(n)}
  +
  F_{\theta}^{(n)}(\vz_{1:L}^{(n)}),
\]
where $F_{\theta}^{(n)}$ contains the temporal mixing operation plus channel mixing and
nonlinearities. This residual structure is not unique to CNNs, but it is essential for
deep sequence models because it stabilizes optimization and lets layers refine rather
than completely replace representations.

\section{CNNs, RNNs, And Transformers}

It is useful to compare the three familiar sequence-modeling patterns before introducing
SSMs.

\begin{table}[h]
\centering
\begin{tabular}{llll}
\toprule
Model family & Temporal mixing & Parallel over time? & Context pattern \\
\midrule
CNN & fixed kernel & yes & local or dilated \\
RNN & hidden recurrence & no, in ordinary form & compressed past state \\
Transformer & attention weights & yes, during training & direct pairwise access \\
\bottomrule
\end{tabular}
\caption{A simplified comparison of common sequence-modeling mechanisms.}
\label{tab:cnn-rnn-transformer}
\end{table}

CNNs are efficient because the same local filter is applied everywhere. RNNs are naturally
online because they update a state one step at a time. Transformers are flexible because
attention weights depend on the content of the sequence, but causal self-attention over
length $L$ has an $L \times L$ interaction pattern.

Deep SSMs occupy a middle ground. Like RNNs, they maintain a state and can be run online.
Like CNNs, linear time-invariant SSMs can be expressed as convolutions and trained in
parallel. Like neither plain CNNs nor plain RNNs, modern SSMs use carefully structured
state matrices to produce long, stable filters.

\section{The Bridge To State Space Models}

The most important equation in this chapter is still the causal convolution:
\[
  \vy_t = \sum_{i=0}^{K-1} \mK_i \vu_{t-i}.
\]
For an ordinary CNN, the matrices $\mK_i$ are learned directly as free parameters. If
$K=7$, the model learns seven kernel slices. If $K=1024$, direct learning becomes much
heavier.

A linear state space model instead defines a recurrence
\[
  \vx_t = \mA \vx_{t-1} + \mB \vu_t,
  \qquad
  \vy_t = \mC \vx_t + \mD \vu_t.
\]
If we unroll this recurrence, we get a convolution-like expression:
\[
  \vy_t
  =
  \mD\vu_t
  +
  \sum_{i=0}^{t-1} \mC \mA^i \mB \vu_{t-i}.
\]
The kernel is no longer directly learned as arbitrary matrices. It is generated by
\[
  \mK_i = \mC \mA^i \mB.
\]
This is the central bridge. A CNN says, ``learn the filter.'' An SSM says, ``learn a
dynamical system whose impulse response is the filter.'' S4, S5, and Mamba differ in how
they parameterize this dynamical system, how they compute it efficiently, and whether the
parameters are fixed across time or depend on the input.

\section{What To Remember}

\begin{itemize}
  \item A temporal convolution is a shared filter applied across sequence positions.
  \item Causal convolution only uses present and past inputs.
  \item Kernel size, stride, dilation, and depth determine the receptive field.
  \item Full, depthwise, grouped, and pointwise convolutions make different tradeoffs
  between temporal mixing and channel mixing.
  \item Convolution can be represented as multiplication by a structured Toeplitz matrix.
  \item A finite convolution kernel is an impulse response.
  \item SSMs generate convolution kernels from recurrence parameters, which is why this
  chapter is the right starting point for S4 and S5.
\end{itemize}

\section{Exercises}

\begin{enumerate}
  \item For a stack of four causal convolutions with kernel size $5$, stride $1$, and no
  dilation, compute the receptive field size.
  \item Write the Toeplitz matrix for a scalar causal convolution with kernel
  $(k_0,k_1,k_2,k_3)$ applied to a length-$5$ input sequence with left zero padding.
  \item Suppose neural data are binned at $20$ ms. What time span is covered by a causal
  convolution with kernel size $9$? What changes if the dilation is $2$?
  \item Explain in words why a causal convolution is parallel during training but still
  respects online decoding constraints.
  \item Compare a learned CNN kernel $\mK_i$ with an SSM-generated kernel
  $\mC\mA^i\mB$. What is gained and what is constrained by the SSM parameterization?
\end{enumerate}
